\section{Problem Set 5 - 15 March 2026}
\subsection{The Quantum Fourier Transform Circuit}
This problem explores the circuit for the Quantum Fourier Transform (QFT). \\

For $N = 2^n$, the QFT can be defined recursively. If we separate the input state $\ket{x_1x_2\ldots x_n}$ into the first bit and the remaining bits, the transformation can be expressed as:
\begin{equation}
    QFT_n\ket{x_1x_2\ldots x_n} = \frac{1}{\sqrt{2}}\left(\ket{0} + e^{2\pi i [0.x_1x_2\ldots x_n]}\ket{1} \otimes QFT_{n-1}\ket{x_2\ldots x_n}\right)
\end{equation}
 
In this recursive structure, the first qubit undergoes a Hadamard transform followed by a series of controlled phase shifts $R_k$, where each shift depends on the values of the subsequent qubits $x_2, \ldots, x_n$. Specifically, a rotation $R_k$ is applied to the first qubit, controlled by the $k$-th qubit. This structure repeats for each subsequent qubit.

\subsubsection{QFT Circuit Implementation}
The QFT circuit uses Hadamard gates ($H$) and controlled-phase rotations ($R_k$). The rotation gate is defined as
\begin{equation}
    R_k = \begin{pmatrix}
        1 & 0 \\ 0 & e^{2\pi i/2^k}
    \end{pmatrix}
\end{equation}

Which of the following statements about the QFT circuit are true? \\
\noindent\fbox{%
  \parbox{\dimexpr\linewidth-2\fboxsep-2\fboxrule}{%
    \begin{itemize}
        \item The circuit requires $O(n^2)$ gates for $n$ qubits. \textbf{TRUE}
        \item The circuit naturally outputs the bits in reversed order (LSB to MSB). \textbf{TRUE}
        \item A Hadamard gate is applied to every qubit exactly once. \textbf{TRUE}
        \item The controlled-$R_k$ gates are performed after all the Hadamard operations. \textbf{FALSE} 
    \end{itemize}
    }%
} 
\newline \\ \\
In cirq, the QFT is implemented by iterating through each qubit, applying a Hadamard, and then iterating through the remaining qubits to apply the conditional rotations. Note that the "distance" between the control and target qubits determines the power of the rotation. \\

Complete the function below to return a Cirq circuit implementing an $n$-qubit QFT on the provided qubits. Do not include the final swap gates.


\begin{lstlisting}
import cirq
import numpy as np

def qft_no_swaps(nqubits):
    circuit = cirq.Circuit()
    qubits = cirq.LineQubit.range(nqubits)
    for i in range(nqubits):
        # Apply Hadamard to the current qubit
        circuit.append(cirq.H(qubits[i]))
        # Write your code here
        for j in range(i + 1, nqubits):
            exponent = 1 / (2 ** (j - i))
            circuit.append(cirq.CZPowGate(exponent=exponent)(qubits[j], qubits[i]))
    return circuit

\end{lstlisting}


\newpage
\subsection{Spectral Analysis of the QFT Operator}
The Quantum Fourier Transform (QFT) operator $\mathcal{F}$ acting on an $N$-dimensional Hilbert space is defined by its matrix elements $F_{jk} = \frac{1}{\sqrt{N}}\omega^{jk}$, where $\omega = e^{2\pi i /N}$.

\subsubsection{The Characteristic Equation}
Let $N=4$. What is the state $\ket{\psi} = \mathcal{F}^2 \ket{3}$. \\

\noindent\fbox{%
  \parbox{\dimexpr\linewidth-2\fboxsep-2\fboxrule}{%
    Calculating $\mathcal{F}^2$,
    \begin{equation}
        (\mathcal{F}^2)_{jk} = \sum_{l=0}^{N-1} \left(\frac{1}{\sqrt{N}} \omega^{jl}\right) \left(\frac{1}{\sqrt{N}} \omega^{lk}\right) = \frac{1}{N}\sum_{l=0}^{N-1} \omega^{jl}\omega^{lk} =\frac{1}{N}\sum_{l=0}^{N-1} \omega^{l(j+k)} = \frac{1}{N}\sum_{l=0}^{N-1} e^{2\pi i l(j+k)/N} 
    \end{equation}
    Let $m = j+k$, where $m \mod N \not\equiv 0$, then
    \begin{equation}
        (\mathcal{F}^2)_{jk} = \frac{1}{N}\sum_{l=0}^{N-1} e^{2\pi i l m /N} = \frac{1}{N} \frac{1 - e^{2\pi i m}}{1 - e^{2\pi i m /N}} = 0  
    \end{equation}
    Let $m = j+k$, where $m \mod N \equiv 0$, then writing $m=Nq$, where $q$ is an integer
    \begin{equation}
        (\mathcal{F}^2)_{jk} = \frac{1}{N}\sum_{l=0}^{N-1} e^{2\pi i l (Nq) /N} = \frac{1}{N} \sum_{l=0}^{N-1} e^{2\pi i l q}= \frac{1}{N} \sum_{l=0}^{N-1} 1
    \end{equation}
    Therefore, to find $\ket{\psi} = \mathcal{F}^2 \ket{3}$, we must find $ m \mod N  = j+k \mod N = 0$. We are given that $j = 3$, $N = 4$. Therefore, $(j+k) \mod N = 3 + k \mod N = 0$. Therefore $k = 1$. Therefore
    \begin{equation}
        \ket{\psi} = \mathcal{F}^2 \ket{3} = \ket{1}
    \end{equation}
    
    
    Useful fact: Standard root of unity: Let $\omega = e^{2 \pi i / N}$, then
    \begin{equation}
        \sum_{k=0}^{N-1}\omega^{km} = 
        \begin{cases}
            N, & m \mod N \equiv 0 \\
            0, & m \not\equiv 0 
        \end{cases}
    \end{equation}
    }%
} 
\newline \\\\
What is the state $\ket{\psi} = \mathcal{F}^4\ket{3}$? \\
\noindent\fbox{%
  \parbox{\dimexpr\linewidth-2\fboxsep-2\fboxrule}{%
        \begin{align}
            \mathcal{F}^4 &= \sum_{l,m,n=0}^{N-1}\left(\frac{1}{\sqrt{N}} \omega^{jl}\right) \left(\frac{1}{\sqrt{N}} \omega^{lm}\right)\left(\frac{1}{\sqrt{N}} \omega^{mn}\right) \left(\frac{1}{\sqrt{N}} \omega^{nk}\right) \\
            &=\frac{1}{N^2} \sum_{l,m,n=0}\omega^{l(j+m)}\omega^{n(m+k)} \\
            &=\frac{1}{N^2} \sum_{m=0}^{N-1}\left(\sum_{l=0}^{N-1}\omega^{l(j+m)}\right)\left(\sum_{n=0}^{N-1}\omega^{n(m+k)}\right)
        \end{align}
        Therefore, $j+m \equiv 0 \mod N$ and $m+k \equiv 0 \mod N$. In this case, $j = 3$ and $N = 4$ are given, therefore, $m=1$ and $k=3$.  \\

        Alternatively, $\mathcal{F}^2 \ket{j} = \ket{-j \mod N}$, since from standard root of unity, $j + k \equiv 0 \mod N \Rightarrow k = -j \mod N$. \\
        Therefore, $\mathcal{F}^4 \ket{j} = (\mathcal{F}^2)^2 = \ket{-(-j) \mod N} = \ket{j \mod N}$, Therefore, $\mathcal{F}^4 = I$.
    }%
} 
\newline \\\\


Based on what you have just worked out, what are the possible eigenvalues of the QFT operator? \\
\noindent\fbox{%
  \parbox{\dimexpr\linewidth-2\fboxsep-2\fboxrule}{%
        From above, we can see that the eigenvalue of $\mathcal{F}^2$ is 1 and $\mathcal{F}^4$ is 1. Therefore all values $x$ that satisfy $\sqrt{x} = 1$ and $\sqrt[4]{x} = 1$. Therefore, the values are $1, -1, i, -i$. 
    }%
} 



\subsubsection{Multiplicity of Eigenvalues}
The trace of the QFT matrix, $\Tr(\mathcal{F})$, is known as a \textbf{Gauss Sum}. For a given $N$, the trace is given by:
\begin{equation}
    S(N) = \Tr(\mathcal{F}) = \frac{1}{\sqrt{N}} \sum_{j=0}^{N-1}e^{2 \pi i j^2 / N }
\end{equation} 
For the specific case of $N=4$, calculate the numerical value of the Gauss sum $S(N)$:\\
\noindent\fbox{%
  \parbox{\dimexpr\linewidth-2\fboxsep-2\fboxrule}{%
        \begin{align}
             S(4) &= \Tr(\mathcal{F}) = \frac{1}{\sqrt{N}} \sum_{j=0}^{N-1}e^{2 \pi i j^2 / N } \\
             &= \frac{1}{2} \left(e^{2\pi i 0^2 / 4} + e^{2\pi i 1^2 / 4} + e^{2\pi i 2^2 / 4} + e^{2\pi i 3^2 / 4}\right) \\
             &= \frac{1}{2} \left(1 + i + 1 + i\right)\\
             S(4)&=1+i
        \end{align}
    }%
} 

\subsubsection{Constructing Eigenstates}
Quantum states that are invariant under the QFT are known as "Fourier-invariant" states. \\

Let $\ket{\psi}$ be an arbitrary state. Show that the state:
\begin{equation}
    \ket{\phi} = \ket{\psi} + \mathcal{F}\ket{\psi} + \mathcal{F}^2\ket{\psi} + \mathcal{F}^3\ket{\psi} 
\end{equation}

is an eigenstate of the QFT. What is the corresponding eigenvalue? \\
\noindent\fbox{%
  \parbox{\dimexpr\linewidth-2\fboxsep-2\fboxrule}{%
        We want to show that $\mathcal{F}\ket{\phi}$ is an eigenstate of $\mathcal{F}$, i.e., $F\ket{\phi} = \lambda \ket{\phi}$ for some eigenvalue $\lambda$.
        \begin{align}
            \mathcal{F}\ket{\phi} & = \mathcal{F}(\ket{\psi} + \mathcal{F}\ket{\psi} + \mathcal{F}^2\ket{\psi} + \mathcal{F}^3\ket{\psi} ) \\ 
            &= \mathcal{F}\ket{\psi} + \mathcal{F}^2\ket{\psi} + \mathcal{F}^3\ket{\psi} + \mathcal{F}^4\ket{\psi} \\
            &= \mathcal{F}\ket{\psi} + \mathcal{F}^2\ket{\psi} + \mathcal{F}^3\ket{\psi} + \ket{\psi} \\
            &=\ket{\psi} + \mathcal{F}\ket{\psi} + \mathcal{F}^2\ket{\psi} + \mathcal{F}^3\ket{\psi}  \\
            \mathcal{F}\ket{\phi}&=\ket{\phi}
        \end{align}
        The corresponding eigenvalue is 1.
    }%
} 
\newline \\\\

If we choose $\ket{\psi} = \ket{00}$ in a 2-qubit system ($N=4$), analytically compute the resulting normalized eigenstate $\ket{\phi}$. \\
\noindent\fbox{%
  \parbox{\dimexpr\linewidth-2\fboxsep-2\fboxrule}{%
        In this case, $\ket{\psi} = \ket{00} = \ket{0}$. Then,
        \begin{align}
            \ket{\phi} & = \ket{\psi} + \mathcal{F}\ket{\psi} + \mathcal{F}^2\ket{\psi} + \mathcal{F}^3\ket{\psi} \\
            &= \ket{0} + \mathcal{F}\ket{0} + \mathcal{F}^2\ket{0} + \mathcal{F}(\mathcal{F}^2\ket{0}) \\
            &= \ket{0} + \mathcal{F}\ket{0} + \mathcal{F}^2\ket{0} + \mathcal{F}\ket{0} \\
            &= \ket{0} + \frac{1}{2}(\ket{0} + \ket{1} + \ket{2} + \ket{3}) + \ket{0} + \frac{1}{2}(\ket{0} + \ket{1} + \ket{2} + \ket{3}) \\
            &= 3\ket{0} + (\ket{1} + \ket{2} + \ket{3})
        \end{align}
        The normalization is given by $\norm{\braket{\phi|\phi}} = 3^2 + 1^2 + 1^2 + 1^2 = 12$.\\
        Therefore,
        \begin{equation}
            \ket{\phi}  = \frac{1}{\sqrt{12}} (3\ket{00} + \ket{01} + \ket{10} + \ket{11})
        \end{equation}
    }%
} 



\newpage
\subsection{Quantum Arithmetic in the Fourier Domain}
The \textbf{Draper Adder} is an algorithm that performs addition in the ``phase" or Fourier domain. Unlike classical ripple-carry adders, it does not require auxiliary (ancilla) qubits for carry bits, making it highly space-efficient for quantum hardware. \\ 

\textbf{Background and Notation}\\
Consider two $n$-bit integers, $a$ and $b$. We represent these in two quantum registers $\ket{a}$ and $\ket{b}$, where:
\begin{itemize}
    \item $\ket{a} = \ket{a_{n-1}a_{n-2}\ldots a_0}$
    \item $\ket{b} = \ket{b_{n-1}b_{n-2}\ldots b_0}$
\end{itemize}
We first apply the Quantum Fourier Transform (QFT) to the second register to obtain the state $\ket{\phi(b)} = \mathrm{QFT}\ket{b}$. The $l$-th qubit of the register $\ket{\phi(b)}$ (where $l \in \{1,\ldots, n\}$ is in the state:
\begin{equation}
    \frac{1}{\sqrt{2}} (\ket{0} + e^{2 \pi i (b/2^l)}\ket{1})
\end{equation}
 
Note that $b/2^l$ can be written in binary fraction notation as $0.b_{l-1}b_{l-2}\ldots b_0$.\\

\textbf{The Task}\\
To perform addition, we want to evolve the state of the $l$-th qubit of the second register such that its phase reflects the sum $(a+b)$. Specifically, we want the final phase of the $l$-th qubit to be:
\begin{equation}
    \exp\left(2\pi i \frac{a+b}{2^l}\right)
\end{equation}
 
After applying this phase shift, we apply an inverse QFT to obtain:
\begin{equation}
    \mathrm{QFT}^{-1}\ket{\phi(a+b)} = \ket{(a+b) \mod 2^n}
\end{equation}

The final state of the second register is the binary representation of the sum $a+b$, while the first register $\ket{a}$ remains unchanged. Because this method uses the phase of the qubits to ``store" the carry information implicitly, it requires \textbf{zero ancilla qubits}, making it notably more space-efficient than classical-style ripple-carry adders on quantum hardware. \\

\textbf{Analytical Derivation}: Show that by applying a sequence of controlled-phase rotations $R_k$ (where $R_k$ applies a phase $e^{2\pi i / 2^k}$) to the $l$-th qubit of register $b$, controlled by the $j$-th bit of register $a$, we can achieve the desired phase. Determine the $k$ in terms of $l$ and $j$ to give the desired result: \\
    \noindent\fbox{%
      \parbox{\dimexpr\linewidth-2\fboxsep-2\fboxrule}{%
            $R_k$ is given by 
            \begin{equation}
                R_k = \begin{pmatrix}
                    1 & 0 \\ 0 & e^{2\pi i/2^k}
                \end{pmatrix}
            \end{equation}
            The $l$-th qubit of the register $\ket{\phi(b)}$ (where $l \in \{1,\ldots, n\}$ is in the state:
            \begin{equation}
                \frac{1}{\sqrt{2}} (\ket{0} + e^{2 \pi i (b/2^l)}\ket{1})
            \end{equation}
            If we apply $R_k$ controlled by the $j$-th bit of register $a$, then if $R_k\ket{0} = \ket{0}, R_k\ket{1}= e^{2\pi i/2^k} \ket{1}$, succinctly the phase is given by $e^{2\pi i a_j/2^k}$
            \begin{equation}
                \frac{1}{\sqrt{2}} (\ket{0} + e^{2 \pi i (b/2^l)} e^{2\pi i a_j/2^k}\ket{1})
            \end{equation}
            If we let $k = l-j$, then 
            \begin{align}
                \frac{1}{\sqrt{2}} (\ket{0} + e^{2 \pi i (b/2^l)} e^{2\pi ia_j/2^{l-j}}\ket{1}) = \frac{1}{\sqrt{2}} (\ket{0} + e^{2 \pi i (b+ a_j2^j)/2^l} \ket{1})
            \end{align}
            Multiplying over all $j < l$
            \begin{equation}
                \frac{1}{\sqrt{2}} (\ket{0} + e^{2 \pi i (b+ a)/2^l} \ket{1})
            \end{equation}\\

            Review:\\
            Split bits into $j < l$ and $j \leq l$. So
            \begin{equation}
                \frac{a}{2^l} = \sum_{j=0}^{l-1}a_j2^{j-l} + \sum_{j=l}^{n+1}a_j^{j-1}
            \end{equation}
            For $j < l$,
            \begin{equation}
                2^{j-l} = \frac{1}{2^{l-j}}
            \end{equation}
            These terms are the fractional part.
            For $j < l$,
            \begin{equation}
                2^{j-l}
            \end{equation}
            This is an integer. We only care about the $j < l$ parts since $e^{2 \pi i \cdot \text{integer}}= 1$. 
        }%
    } 
    \newline \\\\
\textbf{Cirq Implementation}: Write a Python function using cirq that implements a 2-qubit Draper Adder, by completing the code below: \\
    \begin{lstlisting}
        import cirq

        def draper_adder_2qubit(a_val, b_val):
            # Qubits 0,1: Register A | Qubits 2,3: Register B
            qubits = cirq.LineQubit.range(4)    
            # a_val, b_val should be in range [0, 3]
            a, b = qubits[:2], qubits[2:]       
            circuit = cirq.Circuit()
        
            # 1. QFT on register b (Simplified for 2 qubits)
            circuit.append([
                cirq.H(b[0]), 
                cirq.CZ(b[1], b[0])**0.5, 
                cirq.H(b[1])
            ])
        
            # 2. Phase Addition: Apply rotations to b based on a
            for i in range(2):
              for j in range(2):
                      exponent = 1 / (2 ** (j - i))
                      circuit.append(cirq.CZPowGate(exponent=exponent)(b[i], a[j]))
        
        
            # 3. Inverse QFT on register b
            circuit.append([
                cirq.H(b[1]), 
                cirq.CZ(b[1], b[0])**-0.5, 
                cirq.H(b[0])
            ])

            # 4. Measure the result register
            circuit.append(cirq.measure(*b, key='sum'))  
            # 5. Simulate starting from state |a_val><b_val|
            initial_state = (a_val << 2) | b_val         
            result = cirq.Simulator().simulate(circuit, initial_state=initial_state)
            
            # Extract the measurement from the result object
            # Returns the integer value of the bits measured in register 'b'
            return result.measurements['sum'][0] * 2 + result.measurements['sum'][1]

    \end{lstlisting}











\newpage
\subsection{Shor's Algorithm as Unitary Phase Estimation}
In this problem, we re-examine Shor’s factoring algorithm by treating the modular multiplication step as a phase estimation task. Let $N$ be the integer we wish to factor, and let $a$ be a coprime integer ($\gcd(a,N) = 1$). \\

We define the unitary operator $U_a$ acting on an $L$-qubit register ($L \approx \log_2 N$) as: 
\begin{equation}
    U_a\ket{x} = \ket{ax \mod N}
\end{equation}

\subsubsection{Eigenstates of the Modular Multiplication Operator}
The operator $U_a$ permutes the elements of the group $\mathbb{Z}_N^*$. Let $r$ be the order of $a$ modulo $N$ (i.e., $a^r \equiv 1 \mod N$).

\begin{enumerate}
    \item For each $s \in \{0, 1, \ldots, r-1\}$, define the state:
    \begin{equation}
        \ket{u_s} = \frac{1}{\sqrt{r}}\sum_{k=0}^{r-1}\exp\left(\frac{-2\pi i s k}{r}\right)\ket{a^k \mod N}
    \end{equation}
    Show analytically that $\ket{u_s}$ is an eigenstate of $U_a$. 
    \item Determine the corresponding eigenvalue $\lambda_s$. Express your answer in the form $\lambda_s = e^{2\pi i \phi}$, where $\phi_s$ is the phase we aim to estimate. 
\end{enumerate}
    \noindent\fbox{%
      \parbox{\dimexpr\linewidth-2\fboxsep-2\fboxrule}{%
            \begin{align}
                U_a \ket{u_s} &= \frac{1}{\sqrt{r}}\sum_{k=0}^{r-1}\exp\left(\frac{-2\pi i s k}{r}\right)U_a\ket{a^k \mod N} \\
                &=\frac{1}{\sqrt{r}}\sum_{k=0}^{r-1}\exp\left(\frac{-2\pi i s k}{r}\right)\ket{aa^k \mod N} \\
                &= \frac{1}{\sqrt{r}}\sum_{k=0}^{r-1}\exp\left(\frac{-2\pi i s k}{r}\right)\ket{a^{k+1} \mod N} 
            \end{align}
            Because $r$ is the order of $a \mod N$, the states $\ket{a^0 \mod N}, \ket{a^1 \mod N}, \ldots, \ket{a^{r-1} \mod N}$ are all distinct. Therefore the map $k \mapsto k+1 \mod r$ is just a cyclic relabeling of the same set. 
            \begin{align}
                U_a \ket{u_s} &= \frac{1}{\sqrt{r}}\sum_{k=0}^{r-1}\exp\left(\frac{-2\pi i s (k-1)}{r}\right)\ket{a^k \mod N} \\
                &= \frac{1}{\sqrt{r}}\sum_{k=0}^{r-1}\exp\left(\frac{-2\pi i s k}{r}\right) \exp\left(\frac{2\pi i s}{r}\right)\ket{a^k \mod N}                 
            \end{align}
        } %
    }

\subsubsection{The Superposition of Eigenstates}
In a real implementation of Shor's algorithm, we do not know the period $r$ beforehand, so we cannot easily prepare a specific eigenstate $\ket{u_s}$. Instead, we initialize the target register to $\ket{1}$.
\begin{enumerate}
    \item Prove that the state $\ket{1}$ can be written as a uniform superposition of all $r$ eigenstates:
    \begin{equation}
        \ket{1} = \frac{1}{\sqrt{r}} \sum_{s=0}^{r=1}\ket{u_s}
    \end{equation}

    \item With such an implementation, which of the following statements regarding the measurement of the phase $\phi \approx s/r$ and the extraction of the period $r$ are true? Select all that apply: \\
\end{enumerate}
\noindent\fbox{%
  \parbox{\dimexpr\linewidth-2\fboxsep-2\fboxrule}{%
        \begin{align}
            \frac{1}{\sqrt{r}} \sum_{s=0}^{r-1}\ket{u_s} &= \frac{1}{\sqrt{r}} \sum_{s=0}^{r-1} \left(\frac{1}{\sqrt{r}}\sum_{k=0}^{r-1}\exp\left(\frac{-2\pi i s k}{r}\right)\ket{a^k \mod N}\right) \\
            &= \frac{1}{r} \sum_{k=0}^{r-1}\left(\frac{-2\pi i s k}{r}\ket{a^k \mod N}\right)
        \end{align}
        If $k=0$, then every term is 1, so 
        \begin{equation}
            \sum_{s=0}^{r-1}e^{2\pi i s\cdot 0/r} = \sum_{s=0}^{r-1}1 = r
        \end{equation}
        If $k\neq 0$, then using the geometric series, 
        \begin{align}
            \sum_{s=0}^{r-1}\exp\left(\frac{-2\pi i s k}{r}\right) = \frac{1 - (e^{-2\pi i s k /r})^r}{1 - e^{-2\pi i s k /r}} = 0
        \end{align}
        Therefore, the only term that survives is $k=0$ which means the state is at $\ket{a^0 \mod N} = \ket{1}$. Summing up over $s=0$ to $r-1$, cancels out the $1/r$ therefore,
        \begin{equation}
            \ket{1} = \frac{1}{\sqrt{r}} \sum_{s=0}^{r-1}\ket{u_s}
        \end{equation}
    }%
} \newline \\\\
\noindent\fbox{%
      \parbox{\dimexpr\linewidth-2\fboxsep-2\fboxrule}{%
            \begin{itemize}
                \item The algorithm fails to find $r$ if the measurement collapses the state into an eigenstate where $\gcd(s,r) > 1$. \textbf{FALSE}: \textit{ if $\gcd(s,r) d > 1$, then the fraction is not in lowest terms (for example, 2/6=1/3). If QPE gives $2/6$, continued fractions will recover the reduced fraction $1/3$, whose denominator is $3$ not $6$. But that does not mean the algorithm fails completely.}
                \item If the measurement yields $s=0$, the resulting phase $\phi = 0$ provides no information about $r$ and the trial must be repeated. \textbf{TRUE}:\textit{ if $s=0$, then $\phi = 0/r = 0$. The run says nothing.}
                \item Every value of $s \in \{0,1,\ldots,r-1\}$ is nearly equally likely to be measured with a probability of $1/r$. \textbf{TRUE}: \textit{the amplitude of each eigenstate $\ket{u_s}$ is $1/\sqrt{r}$. So the probability is $1/r$}
                \item The continued fractions algorithm can be used as a post-processing step used to derive the simplified fraction $s/r$ from the QPE output. \textbf{TRUE}.
                \item If $s$ and $r$ are coprime ($\gcd(s,r) = 1$), the denominator of the simplified fraction found via continued fractions is exactly $r$. \textbf{TRUE}.
                \item The state $\ket{1}$ is used because it is an eigenstate of the unitary operator $U$, allowing for a deterministic measurement of $r$. \textbf{FALSE}: \textit{$\ket{1}$ is not an eigenstate of $U$.}
                \item Even if $\gcd(s,r) = d > 1$, we can often recover $r$ by combining results from multiple trials. \textbf{TRUE}.
            \end{itemize}
        }%
    } 
    


\subsubsection{Implementation via Controlled-U Gate}
To perform Phase Estimation, we require the ability to implement controlled versions of $U_{a}^{2j}$
 efficiently. \\

Why is it critical to implement these as modular exponentiations rather than simply applying the base operator $U$, $\ket{x} \rightarrow \ket{ax \mod N}$, repeatedly? Select all that apply: \\
\noindent\fbox{%
  \parbox{\dimexpr\linewidth-2\fboxsep-2\fboxrule}{%
        \begin{itemize}
            \item Applying $U$ repeatedly $2^j$ times would require an exponential number of gates, making the algorithm no faster than classical factoring. \textbf{TRUE}
            \item Modular exponentiation allows us to compute $U^{2^j}$ using only $j$ modular squarings, which keeps the total gate count polynomial in $\log N$. \textbf{TRUE}
            \item Applying $U$ repeatedly is physically unrealistic because quantum gates suffer from decoherence, and a circuit with $2^n$ gates would likely lose its quantum state before finishing. \textbf{TRUE}
            \item Modular exponentiation is only necessary because the phase estimation algorithm requires the input state to be $\ket{1}$ rather than an eigenstate. \textbf{FALSE}
            \item The total number of modular multiplications required for the entire QPE circuit using modular exponentiation is $O(n)$, where $n \approx \log_2 N$. \textbf{TRUE}
            \item Computing $a^{2^j} \mod N$ can be done efficiently by repeatedly squaring the previous result: $a^{2^j} = (a^{2^{}j-1})^2$. \textbf{TRUE}
        \end{itemize}
    }%
} 

\newpage
\subsection{Reversible Computing of a Sequential Function}
Consider the sequential function
\begin{equation}
    f = f_5 \circ f_4 \circ f_3 \circ f_2 \circ f_1  
\end{equation}
that is,
\begin{equation}
    f(x) = f_5(f_4(f_3(f_2(f_1(x)))))
\end{equation}
where each $f_1$  maps classical bit strings to bit strings. In quantum algorithms — for example, in quantum walk algorithms, quantum simulation, and Regev's recent quantum factoring algorithm — we often need to implement $f$  as a unitary (reversible) gate acting on a superposition input. \\

In this problem set you will analyze \textbf{why} garbage-free reversible computation is essential for quantum coherence, \textbf{how} to quantify the space-time cost of doing so using pebbling games, and \textbf{how} quantum ghost pebbling reduces the required space. \\

We begin with why reversible, garbage-free computation is desirable and often necessary for the goals of quantum algorithms.

\subsubsection{Forward Computation and the Garbage Problem}
The standard reversible gate for a function $f_1$ is the controlled mapping
\begin{equation}
    R_{f_1} \! : \! \ket{x}\ket{0} \longmapsto \ket{x}\ket{f_1(x)}
\end{equation}
Starting from five fresh ancilla registers, the \textbf{naive forward computation} of $f$ produces
\begin{equation}
    \ket{x_0}\ket{0}\ket{0}\ket{0}\ket{0}\ket{0} \xrightarrow[]{R_{f_1}, R_{f_2}, R_{f_3}, R_{f_4}, R_{f_5}}\ket{x_0}\ket{x_1}\ket{x_2}\ket{x_3}\ket{x_4}\ket{x_5}
\end{equation}
where $x_k = f_k(x_{k-1})$. The registers $\ket{x_1}$ through $\ket{x_4}$  are intermediate garbage - no longer needed once $\ket{x_5}$ is in hand, but they remain and entangle with the rest of the state. \\

Which of the following correctly describe the state $\ket{x_0}\ket{x_1}\ket{x_2}\ket{x_3}\ket{x_4}\ket{x_5}$? Select all that apply. \\
\noindent\fbox{%
  \parbox{\dimexpr\linewidth-2\fboxsep-2\fboxrule}{%
        \begin{itemize}
            \item The mapping $x_0 \mapsto (x_0, x_1, x_2, x_3, x_4, x_5)$ is injective on classical inputs, so the transformation is reversible. \textbf{TRUE}
            \item If $\ket{x_0} = \frac{1}{\sqrt{2}}(\ket{0} + \ket{1})$, the intermediate registers become entangled with $x_0$ in the resulting state. \textbf{TRUE}
            \item The intermediate registers $\ket{x_1},\ldots,\ket{x_4}$ can be safely traced out (partial-traced) without affecting the coherence between $\ket{x_0}$ and $\ket{x_5}$. \textbf{FALSE}
            \item This forward mapping is a valid unitary to use inside interference-based quantum algorithms without any modification. \textbf{FALSE}
        \end{itemize}
    }%
} 


\subsubsection{The Cost of Ancilla Storage}
How many ancilla registers (qubits) are required by the naive forward computation of $f = f_5 \circ f_4 \circ f_3 \circ f_2 \circ f_1 $? (Count only the ancilla registers, not the input register $\ket{x_0}$.) \\
\noindent\fbox{%
  \parbox{\dimexpr\linewidth-2\fboxsep-2\fboxrule}{%
        The ancilla registers are $\ket{x_1}\ket{x_2}\ket{x_3}\ket{x_4}\ket{x_5}$. Therefore, there are 5 ancilla registers total. 
    }%
} 


\subsubsection{Garbage Destroys Quantum Interference}
Now suppose $\ket{x_0} = \ket{+} = \frac{1}{\sqrt{2}}(\ket{0} + \ket{1})$ and we apply the naive forward computation. We then attempt to exploit quantum interference by applying a Hadamard gate $H$ to register $x_0$. Which of the following correctly describe the outcome? Select all that apply: \\
\noindent\fbox{%
  \parbox{\dimexpr\linewidth-2\fboxsep-2\fboxrule}{%
        \begin{itemize}
            \item The garbage registers $\ket{x_1}, \ldots, \ket{x_4}$ are entangled with $\ket{x_0}$, preventing constructive or destructive interference in $x_0$. \textbf{TRUE}
            \item The probability of measuring $x_0 = 0$ after applying $H$ to $x_0$ is exactly $\frac{1}{2}$ - no useful interference occurs. \textbf{TRUE}
            \item Interference still occurs in $x_0$ because the functions $f_1$ are deterministic, making the state effectively classical. \textbf{FALSE}
            \item The reduced density matrix of $x_0$, after tracing out all other registers, is $\frac{1}{2} I$ - a maximally mixed state. \textbf{TRUE}
            \item If all five $f_i$ happen to be the identity function, garbage still prevents interference. \textbf{TRUE}
        \end{itemize}
    }%
} 





\subsubsection{Uncomputation Restores Coherence}
To use $f$ inside a quantum algorithm we must uncompute the intermediate garbage after copying $\ket{x_5}$ to an output register. The uncomputation proceeds by applying the inverse gates in reverse order:
\begin{equation}
    R_{f_4}^{-1}, R_{f_3}^{-1}, R_{f_2}^{-1}, R_{f_1}^{-1}
\end{equation}
which restores registers $x_1, x_2, x_3, x_4$ to $\ket{0}$ while leaving $\ket{x_5}$ (output) intact. \\

Which of the following correctly describes the complete forward-then-uncompute sequence? Choose the one correct answer: \\
\noindent\fbox{%
  \parbox{\dimexpr\linewidth-2\fboxsep-2\fboxrule}{%
        \begin{itemize}
            \item Apply $R_{f_1}, R_{f_2}, R_{f_3}, R_{f_4}, R_{f_5}$, then $R_{f_5}^{-1}R_{f_4}^{-1}, R_{f_3}^{-1}, R_{f_2}^{-1}, R_{f_1}^{-1}$
            \item Apply $R_{f_1}, R_{f_2}, R_{f_3}, R_{f_4}, R_{f_5}$, then $R_{f_4}^{-1}, R_{f_3}^{-1}, R_{f_2}^{-1}, R_{f_1}^{-1}$, leaving $\ket{x_5}$ in the output register. \textbf{TRUE}
            \item Measure registers $x_1,x_2,x_3,x_4$ in the computational basis and reset them to $\ket{0}$ classically. 
            \item Apply $R_{f_1}^{-1}R_{f_2}^{-1}, R_{f_3}^{-1}, R_{f_4}^{-1}$ before $R_{f_5}^{-1}$ to avoid creating garbage. 
        \end{itemize}
    }%
} \newline \\\\
What is the \textbf{total number of reversible function applications} (i.e., applications of $R_{f_i}$ or $R_{f_i}^{-1}$ for any $i$) in the complete forward-then-uncompute procedure for $n=5$ steps?\\
\noindent\fbox{%
  \parbox{\dimexpr\linewidth-2\fboxsep-2\fboxrule}{%
        For $n=5$, we applied $R_{f_1}, R_{f_2}, R_{f_3}, R_{f_4}, R_{f_5}$, then $R_{f_4}^{-1}, R_{f_3}^{-1}, R_{f_2}^{-1}, R_{f_1}^{-1}$. Therefore, the total function applications = 9.
    }%
} 
\newline \\\\
After the full forward-then-uncompute sequence, the state is
\begin{equation}
    \frac{1}{\sqrt{2}}(\ket{0} + \ket{1}) \otimes \ket{0} \otimes \ket{0} \otimes \ket{0} \otimes \ket{0} \otimes \ket{f(\cdot)} 
\end{equation}
 
Which of the following are true of this final state? Select all that apply: \\
\noindent\fbox{%
  \parbox{\dimexpr\linewidth-2\fboxsep-2\fboxrule}{%
        \begin{itemize}
            \item The state is a product between the input register $x_0$ and the output register $x_5$. \textbf{TRUE}
            \item Applying $H$ to $x_0$ now yields constructive or destructive interference (measurement probabilities 0 or 1). \textbf{TRUE}
            \item The output register $\ket{x_5}$ remains entangled with the garbage registers $\ket{x_1},\cdots,\ket{x_4}$. \textbf{FALSE}
            \item The intermediate ancilla registers are returned to $\ket{0}$ and can be reused in subsequent computations. \textbf{TRUE}
        \end{itemize}
    }%
}

\newpage
\subsection{The Reversible Pebbling Game}
\subsubsection{Setup}
The \textbf{space-time} cost of reversibly computing $f$ is captured by the \textbf{reversible pebble game} on a directed line graph with nodes $\{0,1,2,3,4,5\}$:
\begin{itemize}
    \item Node $0$ represents the \textbf{input} $x_0$ and \textbf{always} has a pebble.
    \item Node $i$ ($i \geq 1$) represents register $x_i = f_i(x_{i-1)}$.
    \item A \textbf{pebble} on node $i$ means that register is occupied (holding its value).
\end{itemize}

Two moves are allowed:
\begin{itemize}
    \item \textbf{Pebble}$(i)$: requires node $i-1$  to have a pebble. Places a pebble on node $i$ (computing $x_i$ from $x_{i-1}$ via $R_{f_i}$).
    \item \textbf{Unpebble}$(i)$: requires both nodes $i-1$ and $i$ to have pebbles. Removes the pebble from node $i$ (uncomputing $x_i$ using $x_{i-1}$ via $R_{f_i}^{-1}$).
\end{itemize}
\textbf{Goal}: Place a pebble on node $5$ at some point (after which we copy $x_5$ to a separate output register), then return to the initial state with \textbf{only node} $0$ \textbf{pebbled}.\\

The \textbf{pebbling cost} is the total number of Pebble $+$ Unpebble operations. The \textbf{pebbling space} is the maximum number of pebbles simultaneously present on nodes $1$ through $5$.


\subsubsection{The Naive (Straight-Line) Strategy}
The naive strategy computes all intermediates forward, copies the output, then uncomputes backward:
\begin{equation}
    \mathrm{Forward: P(1), P(2), P(3), P(4), P(5);} \qquad \mathrm{Backward: UP(4), UP(3), UP(2), UP(1)}
\end{equation}
What is the \textbf{pebbling space} (maximum simultaneous pebbles on nodes $1$–$5$) for the naive strategy? \\
\noindent\fbox{%
  \parbox{\dimexpr\linewidth-2\fboxsep-2\fboxrule}{%
        pebbling space = 5
    }%
} 
\newline \\
What is the \textbf{pebbling cost} (total Pebble $+$ Unpebble operations) for the naive strategy? \\
\noindent\fbox{%
  \parbox{\dimexpr\linewidth-2\fboxsep-2\fboxrule}{%
        pebbling cost = 9
    }%
} 

\subsubsection{Bennett's Checkpoint Strategy}
The naive strategy uses space $S=n$ but only time $T=2n-1$. To reduce space, we can use \textbf{checkpoints}. For $n=5$ with a single checkpoint at node $3$, the strategy is: \\

\textbf{Phase 1} (forward to checkpoint $3$, then clean up):
\begin{equation}
    \mathrm{P(1), P(2), P(3), UP(2), UP(1)}
\end{equation}
Result: only node $3$ is pebbled (the checkpoint). \\

\textbf{Phase 2} (forward from checkpoint to $5$):
\begin{equation}
    \mathrm{P(4), P(5)}
\end{equation}
[Copy $x_5$ to output register.]\\ 

\textbf{Phase 3} (backward from $5$ to checkpoint $3$): \\
\begin{equation}
    \mathrm{UP(5), UP(4)}
\end{equation}

\textbf{Phase 4} (re-forward to checkpoint $3$ and uncompute):
\begin{equation}
    \mathrm{P(1), P(2), P(3), UP(3), UP(2), UP(1)}
\end{equation}
What is the \textbf{total pebbling cost} of Bennett's checkpoint strategy for $n=5$?\\
\noindent\fbox{%
  \parbox{\dimexpr\linewidth-2\fboxsep-2\fboxrule}{%
        pebbling cost = 14
    }%
} 
\newline \\\\
What is the \textbf{pebbling space} (maximum simultaneous pebbles on nodes $1$–$5$) for Bennett's checkpoint strategy?\\
\noindent\fbox{%
  \parbox{\dimexpr\linewidth-2\fboxsep-2\fboxrule}{%
        pebbling space = 3
    }%
} 
\newline \\
Which of the following correctly describe space-time tradeoffs in reversible pebbling? Select all that apply:
\noindent\fbox{%
  \parbox{\dimexpr\linewidth-2\fboxsep-2\fboxrule}{%
        \begin{itemize}
            \item Bennett's single-checkpoint strategy for $n=5$ achieves space $3 < n$ but at a higher time cost ($14>9$ operations) than the naive strategy. \textbf{TRUE}
            \item Using $k$ checkpoints (Bennett's recursive strategy), the time for general $n$ grows as $T = O(n^{1+1/(k-1))}$ for fixed $k$, while space remains $O(k)$. \textbf{TRUE}
            \item There exists a standard (ghost-free) reversible pebbling strategy that achieves both $O(1)$ space and $O(n)$ time simultaneously.\textbf{FALSE}
            \item With space $S = O(\log n)$, Bennett's recursion gives time $T = O(n\cdot \poly \log n)$, far better than $O(n^2)$. \textbf{TRUE}
            \item For general $n$, the naive strategy's time $2n-1$ is optimal (minimum pebbling cost) over all strategies with space $S=n$. \textbf{TRUE}
        \end{itemize}
    }%
}
\newpage
\subsection{Quantum Ghost Pebbling: Measurement-Based Uncomputation}
Standard reversible pebbling computes $f = f_n\circ \cdots \circ f_1$ using $O(n)$ ancilla qubits. A striking alternative is \textbf{measurement-based uncomputation}, introduced by Gidney and developed into spooky pebbling by Kahanamoku-Meyer, Van Kirk, and Ragavan. The core idea is to replace the quantum operation of uncomputing an intermediate register with an $X$-basis measurement that frees the qubit immediately but leaves behind a cheap-to-track classical phase — a \textbf{ghost}. Done carefully, the entire computation of $f(k)$ can be performed coherently and without any long-lived garbage, using only \textbf{3 quantum registers} regardless of $n$, at a cost of $O(n^2)$ total operations.

\subsubsection{The Ghost Measurement}
Suppose we have computed the two-register state
\begin{equation}
    \ket{\Phi} = \sum_k \alpha_k \ket{f_0(k)} \ket{f_1(k)}
\end{equation}
 
where $f_0(k) = k$ is the input and $f_1 = f_1(k)$  is the first intermediate value. We wish to free the $\ket{f_0}$ register without keeping it around as ordinary ancilla garbage. \\

The $X$\textbf{-basis measurement} proceeds as follows:
\begin{enumerate}
    \item Apply Hadamard gates to every qubit of the $\ket{f_0}$ register.
    \item Measure in the computational ($Z$) basis, obtaining a random classical bit string $q_0$.
\end{enumerate}
Using the identity $H\ket{k} = \frac{1}{\sqrt{\abs{K}}} \sum_q (-1)^{k \cdot q} \ket{q}$, the state just before measurement is
\begin{equation}
    \sum_k \alpha_k \left(\frac{1}{\sqrt{\abs{K}}} \sum_{q_0} (-1)^{f_0 \cdot q_0} \ket{q_0} \right) \ket{f_1}
\end{equation}
Measuring $\ket{q_0}$ yields a specific classical outcome $q_0$ (with equal probability over all $q_0$), collapsing to the post-measurement state
\begin{equation}
    \sum_k \alpha_k (-1)^{f_0(k) \cdot q_0} \ket{f_1(k)}
\end{equation}
The $\ket{f_0}$ qubit is now \textbf{free}. Its influence lives on only as the classical phase factor $(-1)^{f_0 \cdot q_0}$, parametrized by the known classical string $q_0$: this is the \textbf{ghost pebble} of $f_0$. \\

Which of the following correctly describe the $X$-basis measurement of the $\ket{f_1}$ register? Select all that apply:
\noindent\fbox{%
  \parbox{\dimexpr\linewidth-2\fboxsep-2\fboxrule}{%
        \begin{itemize}
            \item The measurement outcome $q_0$ is a uniformly random classical bit string - it carries no quantum coherence. \textbf{TRUE}
            \item The measurement irreversibly destroys all information about $f_0(k)$: the phase $(-1)^{f_0\cdot q_0}$ cannot be used to recover it. \textbf{FALSE}
            \item The qubit register that held $\ket{f_0}$ is freed and can be reused for subsequent computation. \textbf{TRUE}
            \item The reduced density matrix of the $\ket{f_1}$ register becomes mixed after the measurement, destroying coherence over $k$. \textbf{FALSE}
        \end{itemize}
    }%
} 
\newline \\\\
The ghost phase $(-1)^{f_j\cdot q_j}$ is the key classical quantity carrying phase information after the $X$-basis measurement. For a single-bit register where $f_j,q_j \in \{0,1\}$, the $\mathbb{F}_2$ inner product reduces to ordinary multiplication: $f_j \cdot q_j = f_jq_j$. When the measurement outcome is $q_j = 1$, the ghost phase simplifies to a function of $f_j$ alone.



\subsubsection{The Forward Sweep: $n=2$ Steps}
Consider $f = f_2 \circ f_1$, i.e., $f(k) = f_2(f_1(k))$, and write $f_0 = k$, $f_1 = f_1(k)$, $f_2 = f(k)$. Starting from $\ket{\Phi_0} = \sum_k \alpha_k \ket{f_0}$, the forward sweep alternates \textbf{extend} (compute the next intermediate using $R_{f_i}$) and $X$\textbf{-measure} (free the previous register) steps:
\begin{align}
    &\sum_{k}\alpha_k \underbrace{\ket{f_0}}_{1 \text{ pebble}} \xrightarrow[]{\text{extend } f_1} \sum_k \alpha_k \underbrace{\ket{f_0}\ket{f_1}}_{2 \text{ pebbles}} \xrightarrow[]{X-\text{meas }f_0,q_0} \sum_k \alpha_k \underbrace{(-1)^{f_0\cdot q_0}}_{\text{ghost}}\underbrace{\ket{f_1}}_{1 \text{ pebble} } \\
    & \notag\\
    & \xrightarrow[]{\text{extend } f_2} \sum_k \alpha_k(-1)^{f_0\cdot q_0} \underbrace{\ket{f_1}\ket{f_2}}_{2 \text{ pebbles}} \xrightarrow[]{X-\text{meas }f_1,q_1} \sum_k \alpha_k \underbrace{(-1)^{f_0 \cdot q_0}(-1)^{f_1 \cdot q_1}}_{2 \text{ ghosts}} \underbrace{\ket{f_2}}_{1 \text{ pebble} }
\end{align}
Defining the \textbf{ghost phase product} $\phi_i(\vec{q}) = \prod_{j=i}^{n-1} (-1)^{f_j \cdot q_j}$, the final state of the forward sweep is compactly
\begin{equation}
    \sum_k \alpha_k \phi_0(\vec{q})\ket{f_2}, \qquad \phi_0(\vec{q}) = (-1)^{f_0\cdot q_0} (-1)^{f_1\cdot q_1}
\end{equation}
 
The \textbf{output register} $\ket{f_2}$ is a full quantum register (a pebble). The two ghost phases $(-1)^{f_0 \cdot q_0}$ and $(-1)^{f_1 \cdot q_1}$ are pure classical data, stored as the known bit strings $q_0$ and $q_1$. \\

What is the intermediate state just \textbf{before} the $X$-measurement of $f_1$ (i.e., after extending to $f_2$ but before measuring)? Choose the one correct answer: \\
\noindent\fbox{%
  \parbox{\dimexpr\linewidth-2\fboxsep-2\fboxrule}{%
        \begin{itemize}
            \item $\sum_k \alpha_k (-1)^{f_0 \cdot q_0} \ket{f_1} \ket{f_2}$. \textbf{TRUE}
            \item $\sum_k \alpha_k \ket{f_1}\ket{f_2}$ - the ghost phase $(-1)^{f_0\cdot q_0}$ has already been erased by this point
            \item $\sum_k \alpha_k (-1)^{f_0 \cdot q_0} (-1)^{f_1 \cdot q_1} \ket{f_2}$ - both ghost phases have already been applied
            \item $\sum_k \alpha_k \ket{f_0} \ket{f_1} \ket{f_2}$ - all three registers are still present
        \end{itemize}
    }%
} 
\newline \\\\
Which of the following correctly describe the state $\sum_k \alpha_k \phi_0(\vec{q})\ket{f_2}$ at the end of the $n=2$ forward sweep? Select all that apply: \\
\noindent\fbox{%
  \parbox{\dimexpr\linewidth-2\fboxsep-2\fboxrule}{%
        \begin{itemize}
            \item Only 1 quantum register (qubit) is occupied - the output $\ket{f_2}$ \textbf{TRUE}
            \item The ghost phase product $\phi_0(\vec{q}) = (-1)^{f_0\cdot q_0}(-1)^{f_1\cdot q_1}$ depends on the specific input $k$ (since $f_0 = k$ and $f_1 = f_1(k)$ bot depend on $k$). \textbf{TRUE}
            \item The classical bit strings $q_0$ and $q_1$ are no longer needed and can be discarded at this point. \textbf{FALSE}
            \item To obtain the clean output $\sum_k \alpha_k \ket{f(k)}$, the ghost phases $\phi_0(\vec{q})$ must still be erased. \textbf{TRUE}
        \end{itemize}
    }%
}

\subsubsection{The Forward Sweep: General $n$ Steps}
For the general $n$-step function $f = f_n \circ \cdots f_1$, the forward sweep iterates the extend+$X$-measure pattern $n$ times:
\begin{equation}
    \sum_k \alpha_k \ket{f_0} \xrightarrow[]{\text{extend, }X\text{-meas, extend, }X\text{-meas,}\cdots}\sum_k \alpha_k \underbrace{\left[ \prod_{j=0}^{n-1}(-1)^{f_j \cdot q_j} \right]}_{\phi_0(\vec{q})\text{, all classical}} \ket{f_n}
\end{equation}
At every step the pattern is: extend (peak at 2 pebbles), then $X$-measure the older register (back to 1 pebble). The result is the output state $\ket{f_n}$ adorned with $n$ classical ghost phases. \\

What is the maximum number of pebbles (live quantum registers) in use at any point during the $n=5$ forward sweep? \\
\noindent\fbox{%
  \parbox{\dimexpr\linewidth-2\fboxsep-2\fboxrule}{%
        max pebbles = 2
    }%
}
\newline \\\\
After the $n$-step forward sweep, the output state carries the accumulated ghost phase product
\begin{equation}
    \phi_0(\vec{q}) = (-1)^{f_0q_0} \cdot (-1)^{f_1q_1} \cdots (-1)^{f_{n-1}q_{n-1}}
\end{equation}
Because $(-1)^a \cdot (-1)^b = (-1)^{a+b}$, this product can be written as a single exponential $\phi_q(\vec{q}) = (-1)^E$. \\

For the $n=2$ case, write the exponent $E$ as an algebraic expression in $f_0, q_0, f_1, q_1$.\\
\noindent\fbox{%
  \parbox{\dimexpr\linewidth-2\fboxsep-2\fboxrule}{%
        \begin{equation}
            E = f_0q_0 + f_1q_1
        \end{equation}
    }%
}

\subsubsection{The Backward Sweep and Phase Erasure}
After the forward sweep we have $\sum_k \alpha_k \phi_0(\vec{q}) \ket{f_n}$ with $n$ ghost phases and 1 pebble. To recover the clean output $\sum_k \alpha_k \ket{f_n}$  we must erase all ghost phases. The \textbf{backward sweep} does this one ghost at a time:\\

\textbf{One iteration of the backward sweep} (erasing the $f_0$ ghost, illustrated for general $n$):
\begin{enumerate}
    \item \textbf{Reverse to} $f_1$: Starting from $\ket{f_n}$, apply the inverse circuits in sequence, keeping $\ket{f_n}$ and adding new registers while $X$-measuring each one as it is produced. This traverses back through $f_{n-1},f_{n-2},\ldots,f_1$, converting each along the way into a new ghost phase. The peak state is $f_{n-1}\ket{f_{j-1}} \ket{f_j}$ - \textbf{3 pebbles} - which briefly appears as each new $\ket{f_{j-1}}$ is computed from $\ket{f_j}$ before $\ket{f_j}$ is measured.
    \item \textbf{Phase erasure}: With $\ket{f_n}$ and $\ket{f_1}$ simultaneously present, re-compute $f_0 = f_1^{-1}(f_1)$ classically (using the inverse of $f_1$) and apply the phase correction $(-1)^{-f_0\cdot q_0}$  as a classically controlled phase gate. This cancels the $(-1)^{f_0q_0}$ ghost.
    \item \textbf{Measure} $f_1$: $X$-measure $\ket{f_1}$ with outcome $q_1'$, producing a new ghost phase $(-1)^{f_1\cdot q_1 '}$ and freeing the $\ket{f_1}$ register. Update the classical record: $q_1 \leftarrow q_1 + q_1'$.
\end{enumerate}
After one iteration: the $f_0$ ghost is erased, the remaining $n-1$ ghost phases are updated, and we are back to 1 pebble. Repeating $n$ times in total erases all ghosts and yields the clean output $\sum_k \alpha_k \ket{f_n}$.\\


What is the \textbf{maximum number of pebbles} simultaneously in use during the backward sweep? Choose the one correct answer: \\
\noindent\fbox{%
  \parbox{\dimexpr\linewidth-2\fboxsep-2\fboxrule}{%
        \begin{itemize}
            \item 1 - only $\ket{f_n}$ is ever present during the backward sweep.
            \item 2 - the backward sweep reuses the same 2-pebble pattern as the forward sweep.
            \item 3 - during each ``reverse" step, $\ket{f_n}$, $\ket{f_j-1}$, and $\ket{f_j}$ are all live simultaneously. \textbf{TRUE}
            \item $n$ - all intermediate values must be reconstructed at once before erasing the ghost. 
        \end{itemize}
    }%
}
\newline \\\\
What is the maximum number of pebbles in use at any point during the entire ghost pebbling protocol (forward sweep + all backward sweep iterations) for $n=5$? \\
\noindent\fbox{%
  \parbox{\dimexpr\linewidth-2\fboxsep-2\fboxrule}{%
        max pebbles over entire protocol = 3
    }%
}


\subsubsection{Resource Costs and Comparison}
We can now compare the resource costs of the three main strategies for reversibly computing $f = f_n \circ \cdots \circ f_1$  and returning to a clean state. The total operations count all applications of $R_{f_i}$ or $R_{f_i}^{-1}$ (plus $X$-measurements, which are inexpensive).
\begin{table}[!h]
\begin{tabular}{lcc}
\textbf{Strategy}                      & \multicolumn{1}{l}{\textbf{Max qubits (space)}} & \multicolumn{1}{l}{\textbf{Total operations (time)}} \\
Naive (store all intermediates)        & $n$                                             & $2n-1$                                               \\
Bennett recursive (divide-and-conquer) & $O(\log n)$                                     & $O(n^{1.58})$                                         \\
Ghost pebbling (this section)          & \textbf{3}                                      & $O(n^2)$                                            
\end{tabular}
\end{table}

Ghost pebbling achieves the best space (only 3 qubits, constant in $n$), at the cost of $O(n^2)$ total operations - worse than Bennett's $O(n^{1.58})$ but still polynomial. The key practical advantage is that classical bit manipulation (storing and updating $\vec{q}$) is vastly cheaper than storing $O(\log n)$ live quantum registers. \\

Which of the following correctly summarize the trade-offs among the three strategies for $n=5$? Select all that apply: \\
\noindent\fbox{%
  \parbox{\dimexpr\linewidth-2\fboxsep-2\fboxrule}{%
        \begin{itemize}
            \item Ghost pebbling uses fewer qubits than both the naive strategy and Bennett's strategy for $n=5$. \textbf{FALSE}
            \item The ghost pebbling forward sweep (only) requires fewer operations than the naive strategy ($5$ vs. $9$ extend ops). \textbf{TRUE}
            \item Ghost pebbling's $O(n^2)$ total operation count makes it slower than the naive strategy's $2n-1$ for small $n$, but its space advantage ($3$ qubits vs. $5$) can still be decisive for large $n$. \textbf{TRUE}
            \item Ghost pebbling and Bennett's strategy both use $3$ qubits for $n=5$; the difference is that ghost pebbling replaces live qubit storage with cheap classical bit storage for large $n$. \textbf{TRUE}
            \item The ghost phase $(-1)^{f_j \cdot q_j}$ requires storing the full quantum state of $f_j$ as classical data. \textbf{FALSE}
        \end{itemize}
    }%
}




\newpage
\subsection{Factoring via Lattice Reduction with Partial Bit Knowledge}
In this problem we factor $N=91$ using the quantum period-finding approach from lecture, aided by partial knowledge of the prime factors. The key computational step is the \textbf{extended Euclidean algorithm (EEA)}, which is precisely a two-dimensional lattice reduction: it finds the shortest lattice vector in the span of $(M,0)$ and $(Y,1)$ (where $M$ is the Hilbert-space dimension and $Y$ is the measurement outcome), recovering the hidden period $T$. \\

\textbf{Setup.} We wish to factor $N=91 = p \cdot q$  where $p$ and $q$ are distinct primes. Partial knowledge of the binary representations of $p$ and $q$ is available:
\begin{itemize}
    \item The 4-bit binary representation of $p$ begins with $01$, so $4 \leq p \leq 7$.
    \item The 4-bit binary representation of $q$ begins with $110$, so $12 \leq q \leq 13$.
\end{itemize}
We run the period-finding subroutine on $f(x) = d^x \mod N$ with base $c=2$ and $d \equiv c^2 \equiv 4 \mod 91$.

\subsubsection{Bounding the period}
The period $T$ of $f$ divides $\lambda(N) = \text{lcm}(p-1, q-1)$. The partial bit knowledge gives upper bounds $p \leq p_{\text{max}} = 7$ and $q \leq q_{\text{max}} = 13$, so we can bound
\begin{equation}
    T_{\text{max}} = \text{lcm}(p_{\text{max}} - 1, q_{\text{max}} - 1)
\end{equation}
Compute $T_{\text{max}}$: \\ 
\noindent\fbox{%
  \parbox{\dimexpr\linewidth-2\fboxsep-2\fboxrule}{%
        \begin{align}
            T_{\text{max}} &= \text{lcm}(p_{\text{max}} - 1, q_{\text{max}} - 1)\\ 
            &= \text{lcm}(7 - 1, 13 - 1) \\
            T_{\text{max}} &= \text{lcm}(6, 12)
        \end{align}
        $6 = 2 \times 3$ and $12 = 2^2 \times 3$. Therefore,
        \begin{equation}
            T_{\text{max}} = \text{lcm}(6, 12) = 2^2 \times 3 = 4 \times 3 = 12
        \end{equation}
    }%
}


\subsubsection{Choosing the register dimension}
The inexact quantum period-finding circuit uses a Hilbert space of dimension $M = 4T_{\text{max}^2}$. This guarantees that the Fourier-peak width $\sigma' = M / (2 \pi \sigma)$ is less than $1$ when $\sigma = T_{\text{max}^2}$, so the measured integer $Y$ satisfies
\begin{equation}
    \left| Y - k' \frac{M}{T} \right| < 2 \Rightarrow \left| \frac{Y}{M} -  \frac{k'}{T} \right| < \frac{1}{2T_{\text{max}}^2}
\end{equation}
\noindent\fbox{%
  \parbox{\dimexpr\linewidth-2\fboxsep-2\fboxrule}{%
        \begin{equation}
            M = 4T_{\text{max}}^2
        \end{equation}
    }%
}
\subsubsection{Recovering the period via the EEA}
The quantum circuit with $M = 577$ (a prime slightly above $4T_{\text{max}}=576$) measures the outcome $Y = 481$. . By the Diophantine-approximation guarantee, the fraction $k'/T$ is a \textbf{convergent} of $Y/m$.\\

We find convergents by running the EEA, tracking integer linear combinations $a_iM+b_iY=r_i$  (each row is a lattice vector of decreasing length; the fifth column shows the convergent $k'/T \approx -a_i/b_i$:

\begin{center}
\begin{table}[!h]
\begin{tabular}{ccccc}
\textbf{Step} & $r_i$ & $a_i$ & $b_i$ & $-a_i/b_i$ \\
0 & 577 & 1 & 0 & -- \\
1 & 481 & 0 & 1 & -- \\
2 & 96 & 1 & -1 & 1/1 \\
3 & \textbf{1} & -5 & \textbf{6} & 5/6  
\end{tabular}
\end{table}
\end{center}

We stop at step 3 because the denominator $6 \leq T_{\text{max}} = 12$, while the next step would give denominator. The stopping rule (from lecture) is: take the convergent with the largest denominator $481 \gg T_{\text{max}}$. \\

The key linear combination at step 3 is
\begin{equation}
    (-5) \cdot 577 + 6 \dot 581 = -2885+2886 = 1
\end{equation}

confirming a residual of $1 \ll M$. What is the period $T$ recovered from the EEA? \\
\noindent\fbox{%
  \parbox{\dimexpr\linewidth-2\fboxsep-2\fboxrule}{%
        The algorithm  says choose the convergent with the largest denominator $T_{\text{max}}=12$. From the table we have that \begin{equation}
            \frac{k'}{T} = \frac{5}{6}
        \end{equation}
        Therefore, $T = 6$.
    }%
}


\subsubsection{Extracting the factors}
With period $T=6$ and $c=2$, form  \begin{equation}
a \equiv c^T \pmod N = 2^6 \bmod 91 = 64.
\end{equation}
One can verify
\begin{equation}
a^2 = 64^2 = 4096 = 45\cdot 91 + 1 \equiv 1 \pmod{91}
\end{equation}
and
\begin{equation}
a = 64 \not\equiv \pm 1 \pmod{91}
\end{equation}

(since $64 \neq 1$ and $64 \neq 90$). \\

Therefore $N \mid (a-1)(a+1)$ but $N \nmid a\pm 1$, so the prime factors of $N$ split across the two factors and a simple GCD yields a non-trivial divisor:
\begin{equation}
p = \gcd(a-1,N)=\gcd(63,91).
\end{equation}
Compute $\gcd(c^T-1,N)$ with $c=2$, $T=6$, $N=91$: \\
\noindent\fbox{%
  \parbox{\dimexpr\linewidth-2\fboxsep-2\fboxrule}{%
        \begin{align}
            \gcd(c^T - 1,N) &= \gcd(2^6-1,91) \\
            &= \gcd(63,91) 
        \end{align}
        $91 = 7 \times 13$ and $63 = 9 \times 7$. Therefore $\gcd(63,91) = 7$
    }%
}